% =====================================================================
%  TỔNG QUAN BÁO CÁO
% =====================================================================
\chapter*{Tổng quan báo cáo}
\addcontentsline{toc}{chapter}{Tổng quan báo cáo}
\markboth{Tổng quan báo cáo}{Tổng quan báo cáo}

Mô hình ngôn ngữ lớn (LLM) viết trôi chảy nhưng không biết mình đang sai. Hỏi một trợ lý thông thường
\emph{``đang uống metformin thì tối nay uống bia được không''}, nó có thể trả lời bằng một đoạn văn nghe rất hợp
lý mà không nêu được căn cứ nào, hoặc nêu một con số không có trong bất kỳ tài liệu nào. Trong y tế và tài
chính, một câu trả lời như vậy còn nguy hiểm hơn là không trả lời. Tiểu luận này nghiên cứu ba kỹ thuật bổ sung
cho LLM chính phần ``biết'' và ``kiểm chứng'' mà nó thiếu: \textbf{đồ thị tri thức} (Knowledge Graph, KG) lưu tri
thức có cấu trúc và có nguồn; \textbf{sinh tăng cường truy xuất} (Retrieval-Augmented Generation, RAG) đưa văn
bản gốc vào ngữ cảnh trước khi sinh; và \textbf{tác tử thông minh} theo vòng lặp ReAct phối hợp KG, RAG với các
mô hình học sâu chuyên biệt để trả lời câu hỏi nhiều bước.

\section*{Mục tiêu và phạm vi}
Đề bài yêu cầu trình bày khái niệm và kiến trúc của KG, các thế hệ RAG, rồi áp dụng vào hai miền của Việt Nam:
quản lý \textbf{bệnh đái tháo đường típ 2} và phân tích \textbf{rổ cổ phiếu VN30 cùng giá vàng SJC}, cuối cùng
triển khai một AI Agent có giao diện web. Báo cáo này làm trọn chuỗi đó trên máy cá nhân, không gọi một API
LLM bên ngoài nào:

\begin{itemize}
  \item Viết kho tri thức \SL{retrieval_grid/n_docs} tài liệu Markdown (\SL{retrieval_grid/n_chars} ký tự) tổng
        hợp từ Quyết định 5481/QĐ-BYT, ADA, các Nghị định về vàng và báo cáo doanh nghiệp, kèm bộ
        \SL{retrieval_grid/n_questions} cặp hỏi đáp có đáp án nằm nguyên văn trong kho để chấm điểm truy xuất.
  \item Dựng hai đồ thị tri thức có bản thể học: \texttt{VietnamDiabetesKG} (\SL{kg/health/nodes} thực thể,
        \SL{kg/health/edges} quan hệ) và \texttt{VietnamFinanceKG} (\SL{kg/finance/nodes} thực thể,
        \SL{kg/finance/edges} quan hệ), xuất RDF, suy diễn RDFS và truy vấn SPARQL.
  \item So sánh định lượng 5 chiến lược chia đoạn $\times$ 5 bộ truy xuất, 3 mô hình nhúng, 7 cỡ đoạn,
        hai kỹ thuật tiền truy xuất bằng LLM (HyDE, Multi-Query) và bốn kiến trúc từ Naive RAG tới GraphRAG.
  \item Huấn luyện trên GPU hai mô hình học sâu làm công cụ cho tác tử: MLP chẩn đoán nguy cơ đái tháo đường
        (\SL{diabetes/data/n} hồ sơ) và BiLSTM dự báo lợi suất cho bốn cổ phiếu và vàng quy đổi.
  \item Ghép tất cả thành tác tử ReAct có bảy công cụ, bộ kiểm chứng trích dẫn và con số, đánh giá trên
        \SL{dem/agent_scen} kịch bản, và đóng gói thành Web App Flask.
\end{itemize}

\section*{Phân bổ CPU và GPU}
Máy thực nghiệm có một GPU laptop 8 GB và một CPU 32 luồng. Nguyên tắc phân bổ là: việc nào là phép nhân ma
trận lớn lặp lại nhiều lần thì chạy GPU, việc nào là tra cứu, duyệt đồ thị hay tính toán nhỏ thì để CPU, vì chi
phí chép dữ liệu qua PCIe lớn hơn phần tính toán. Bảng~\ref{tab:moitruong} ghi cụ thể từng thành phần.

\begin{table}[H]\centering
\caption{Môi trường thực nghiệm và phân bổ tính toán}\label{tab:moitruong}
{\small\begin{tabular}{L{4.6cm}L{2.0cm}L{7.6cm}}
\toprule
Thành phần & Thiết bị & Lý do \\
\midrule
Qwen2.5-1.5B-Instruct (bf16) & GPU & sinh văn bản tự hồi quy, mỗi token một lượt qua 28 tầng Transformer \\
bge-reranker-v2-m3 (bf16) & GPU & chấm 30 cặp (câu hỏi, đoạn) mỗi truy vấn bằng cross-encoder 568 triệu tham số \\
multilingual-e5 (fp16 khi chung tiến trình với LLM) & GPU & mã hóa câu hỏi và toàn bộ đoạn khi lập chỉ mục \\
Huấn luyện MLP và BiLSTM & GPU & lô 1024 và 128 mẫu, lan truyền ngược \\
BM25, TF-IDF, chia đoạn & CPU & tính trên chỉ mục thưa, dưới 1 ms mỗi truy vấn \\
Hai KG (NetworkX), RDF, SPARQL, suy diễn & CPU & duyệt đồ thị vài chục nút, rdflib chỉ chạy CPU \\
Suy luận MLP, BiLSTM trong tác tử & GPU & đã nạp sẵn trên GPU cùng LLM; MC Dropout chạy 200 mẫu một lô \\
Vẽ hình, dựng báo cáo & CPU & matplotlib và LuaLaTeX \\
\midrule
\multicolumn{3}{L{14.6cm}}{GPU: NVIDIA GeForce RTX 4060 Laptop 8 GB; CPU: AMD Ryzen 9 7945HX, 32 luồng; RAM
40 GB. Python 3.13, PyTorch 2.13 + CUDA 12.6, sentence-transformers, NetworkX, rdflib, Flask.} \\
\bottomrule
\end{tabular}}
\end{table}

Một ràng buộc thực tế đã định hình kiến trúc: trên Windows, trình điều khiển WDDM tính cả bộ nhớ GPU đã cấp
phát vào hạn mức \emph{commit} của hệ điều hành. Nạp LLM sau embedder và reranker khiến tiến trình vượt hạn mức và
dừng đột ngột khi đọc tệp trọng số. Cách khắc phục được giữ lại trong mã: nạp LLM trước, lưu sẵn reranker ở dạng
bf16 trên đĩa và dùng embedder fp16 khi chung tiến trình với LLM; cả ba mô hình khi đó cùng chạy trong khoảng 8 GB.

\section*{Cấu trúc báo cáo}
Báo cáo bám theo bảy phần của đề bài. \textbf{Chương 1} nêu khái niệm KG, RAG và vai trò của cơ sở tri thức.
\textbf{Chương 2} đi vào cấu tạo KG: thành phần, mô hình bộ ba RDF, bản thể học, suy diễn và so sánh Graph DB
với Vector DB. \textbf{Chương 3} trình bày ba thế hệ RAG, từ Naive RAG, các chiến lược chia đoạn, Advanced RAG tới
GraphRAG, mỗi phần đều có thực nghiệm trên kho tri thức của bài. \textbf{Chương 4} xây hệ thống y tế cho bệnh
đái tháo đường típ 2, \textbf{Chương 5} xây hệ thống tài chính cho VN30 và vàng SJC. \textbf{Chương 6} ghép thành
tác tử ReAct, đánh giá và triển khai Web App. \textbf{Chương 7} kết luận và nêu hướng phát triển.

\begin{phathien}{Quy ước về số liệu và mã nguồn}
Không con số thực nghiệm nào trong báo cáo được gõ tay. Mỗi kịch bản trong \texttt{scripts/} ghi kết quả ra
\texttt{outputs/metrics/*.json}; \texttt{build.py} đọc các tệp đó để sinh macro số liệu, bảng và vết suy luận, rồi
trích nguyên văn từng hàm được trích dẫn từ mã nguồn trong \texttt{src/TieuLuan 02}. Hình dữ liệu vẽ bằng
matplotlib từ cùng các tệp JSON; sơ đồ lý thuyết vẽ bằng TikZ. Chạy lại thực nghiệm rồi dựng lại báo cáo thì mọi
con số, bảng và hình đổi theo.
\end{phathien}
